\begin{abstract}
\textbf{Objective:}
Covariance fitting, Gaussian predictive likelihood, and
unknown-person EEG discrimination are different
regularisation objectives. We examine how a validation
estimand mismatch affects cross-subject conclusions,
and whether correcting or more precisely optimising
the source criterion improves transfer.

\textbf{Methods:}
A source-fitted tangent-space LDA pipeline is evaluated
under participant-grouped, zero-target-information
selection. We audit a historically labelled
covariance-NLL criterion, distinguish its
validation-class-centred score from true
source-mean predictive likelihood, and decompose
their exact difference into common-location,
score-orthogonal, and discriminant-aligned
class-mean shifts. Frozen positive, negative and
post-outcome source-only evidence is kept separate.

\textbf{Results:}
A prospectively frozen EEGMMIDB block (20 people)
showed a $+0.0561$ AUC gain for source-discrimination
selection over the historical class-centred score
(Holm $p=0.00597$), with five materially harmed
participants. In source-validation analyses of
Cho2017, Dreyer2023A and native-montage
BNCI2014\_002, respectively, 98.59\%,
99.29\% and 99.77\% of the corrected
predictive-score mean-shift gap was attributable
to components invariant to the tested fixed
classifier's within-person AUC. These
three mechanism analyses were \emph{post-outcome}.
A ten-person protocol-first Weibo2014 test found
identical held-person AUC for continuous and
27-point source selection (0.73141).
A 40-person Cho2017 sequential follow-up
gave an essentially zero overall source-envelope
gain ($-0.000053$ AUC), not a superior decoder.

\textbf{Conclusion:}
A mathematically proper Gaussian likelihood
can still primarily measure variation irrelevant
to an EEG discriminant's within-person ranking.
Correct estimands, source-selection information
boundaries and individual harms are central;
the tested more elaborate source-only selectors
have not shown stable transfer superiority.
\end{abstract}
